\section*{Bab \ref{ch:g5:people-and-nature} --- Manusia dan Alam}

\begin{solution}{exo:g5:people-and-nature:1}
Ladang, dikelola untuk satu tanaman pilihan; hutan, dikelola untuk
kayu; padang rumput, dijaga tetap terbuka untuk jerami atau
penggembalaan; kota, jalan dan bendungan, dibangun untuk tempat
tinggal dan perjalanan. (Empat mana pun beserta tujuannya.)
\end{solution}

\begin{solution}{exo:g5:people-and-nature:2}
Kalau dibiarkan tanpa dipangkas dan tanpa digembalai, ia tumbuh
kembali menjadi semak belukar dalam beberapa tahun, lalu menjadi hutan
--- \omterm{def:g2:where-animals-live:habitat}{habitat} padang rumput yang terbuka dan berbunga itu lenyap.
\end{solution}

\begin{solution}{exo:g5:people-and-nature:3}
Jangan mengambil lebih cepat daripada pembaruannya; pertahankan para
pekerjanya (\omterm{def:g5:biodiversity:species}{spesies} yang jasa cuma-cumanya menjadi sandaran
pemakaiannya); kembalikan apa yang kamu pinjam (tutuplah gelangnya,
jangan tinggalkan buangan yang menumpuk).
\end{solution}

\begin{solution}{exo:g5:people-and-nature:4}
Pagar tanaman yang dibongkar lawan yang dipertahankan: A kehilangan
penahan tanahnya dan rumah para pemakan hamanya. Tanah musim dingin
yang telanjang lawan yang tertutup: tanah A hanyut menuruni bukit
ketika hujan. Satu tanaman tunggal yang raksasa lawan penggiliran: A
mengundang hama apa pun dari tanaman itu untuk menyapu seluruh
lerengnya, dan tidak pernah mengistirahatkan tanahnya.
\end{solution}

\begin{solution}{exo:g5:people-and-nature:5}
Hitunglah apa yang ditumbuhkan hutannya dalam setahun; tebanglah tidak
lebih daripada itu. Itulah aturan 1 --- jangan mengambil lebih cepat
daripada pembaruannya --- dan hutannya menghasilkan kayu setiap tahun,
tanpa batas waktu, sambil tetap menjadi hutan yang hidup.
\end{solution}

\begin{solution}{exo:g5:people-and-nature:6}
Misalnya: \omterm{prop:g3:sorting-living-things:groups}{burung} walet, yang bersarang di ``tebing'' beton; rubah,
yang berpatroli di tempat sampah dan kebun; \omterm{def:g1:plants-around-us:plant}{tumbuhan} liar, yang
berakar di celah aspal dan sambungan tembok (juga merpati di birai,
elang di menara).
\end{solution}

\begin{solution}{exo:g5:people-and-nature:7}
Bagi manusia: air yang lebih lambat dan penyimpanan banjir di padang
rumput basah yang dipulihkan --- banjir yang lebih sedikit di hilirnya.
Bagi jaringnya: \omterm{def:g2:where-animals-live:habitat}{habitat} keloknya kembali, dan bersamanya kembali pula
\omterm{prop:g3:sorting-living-things:groups}{ikan}, capung serta padang rumput perburuan bangaunya.
\end{solution}

\begin{solution}{exo:g5:people-and-nature:8}
Jawaban terbuka. Temuan yang lazim: \omterm{def:g1:my-body:joint}{bahu} jalan yang dipangkas, pohon
yang dipangkas, petak yang ditanami, trotoar yang disiangi --- dengan
barangkali satu puncak tembok atau parit yang tidak dikelola, yang
sering merupakan tempat paling liar dalam pemandangannya.
\end{solution}

\begin{solution}{exo:g5:people-and-nature:9}
Tidak pernah menyentuh justru akan menghancurkan sebagian \omterm{def:g2:where-animals-live:habitat}{habitat}:
padang rumput yang berbunga itu ada \emph{karena} pemangkasannya, dan
ia menutup tanpa pemangkasan. Melindungi alam bukanlah nol sentuhan
melainkan pengelolaan yang baik --- pertanyaan yang jujur adalah
bagaimana menyentuhnya: aturan mana yang dihormati pemakaiannya,
bukan apakah tangan pernah diletakkan sama sekali.
\end{solution}

\begin{solution}{exo:g5:people-and-nature:10}
Aturan 2 dilanggar terang-terangan (para pekerjanya diusir bersama
pagar tanamannya), aturan 3 sebagian (tanah yang hanyut tidak pernah
dikembalikan), dan aturan 1 terancam (tanahnya hilang lebih cepat
daripada tanah terbentuk). Satu perubahan yang paling memperbaiki:
tanamlah kembali pagar tanamannya --- ia menahan tanahnya dan
memondokkan ulang para pemakan hamanya dalam satu tarikan.
\end{solution}

\begin{solution}{exo:g5:people-and-nature:11}
Tangkapan yang dilipatduakan mengambil \omterm{prop:g3:sorting-living-things:groups}{ikan} jauh lebih cepat daripada
perkembangbiakan gerombolannya --- yang berarti membelanjakan
persediaan pembiakannya sendiri. Membelanjakan modal selalu tampak
gemilang selama ia masih ada: palkanya penuh justru karena \omterm{prop:g3:sorting-living-things:groups}{ikan} tahun
depan sedang didaratkan tahun ini. Sang rimbawan hanya memanen
pertumbuhan tahun itu dan menjaga persediaan tumbuhnya tetap utuh ---
hitung-hitungan yang sama, dengan tanda yang berlawanan, dan itulah
sebabnya hutannya menghasilkan selamanya sedangkan armadanya berhenti
pada tahun keenam.
\end{solution}
